% =============================================================================
% 2.1 Introduction
% =============================================================================
\section{Introduction}
\label{sec:ch2-intro}

La méthodologie adoptée s'inscrit dans une démarche de recherche et
développement appliquée, visant à concevoir, développer et valider une
plateforme de capitalisation et de génération assistée de dossiers d'appels
d'offres. Elle combine des principes du génie logiciel, des méthodes de
l'intelligence artificielle et une analyse du domaine métier des marchés
publics camerounais.

Ce chapitre n'est pas un journal de bord du développement. Il explicite les
choix, leurs justifications et les garanties qu'ils produisent. Un lecteur
compétent doit pouvoir y suivre le problème métier, le flux de traitement
de bout en bout, la place de chaque brique, et la frontière entre le
périmètre de la première version et ce qui est explicitement reporté.

Nous procédons en huit temps~: formulation de la problématique et de
l'approche générale~; analyse des besoins et critères d'évaluation~;
modélisation des données~; architecture du système~; stratégies de
recherche et de génération~; workflow de traitement~; protocole
d'évaluation~; planification.
